\scanpage{9}
où $\delta'\inserted{\in\underline{\mathrm{NS}}_{B/S}(S)}$ a été déjà caractérisé par ailleurs, par
\[
\widetilde{\delta'}\,\widetilde\delta=-\chi(\delta)\,\mathrm{id}_A\qquad\ldots.
\]

On peut se proposer, pour une section $\lambda$ de $\underline{\mathrm{Pic}}_{A/S}$, définissant le faisceau inv. sur $A$, de ``calculer'' la classe ds $K^{\cdot}(B)$ de $\varphi_L(\deleted{\unsure{\delta}}\mathrm{cl}(L))$ en termes de \deleted{\unsure{la se}} \unsure{de|la} section $\varphi(\lambda)^{\unsure{\,\cdot\,|*|\vphantom{*}}}$ de $\mathrm{Pic}_{B/S}$, définissant un faisceau inversible $L'^{-1}$ sur $B$ (de classe $-\delta'$ dans $\underline{\mathrm{NS}}.(B/S)$ si $\lambda$ est de classe $\delta$ dans $\mathrm{NS}(A/S)$). \unsure{Heuristiqmt|Heuristiquement}, on \deleted{\unsure{aura}} \inserted{devrait avoir} (pour $\chi(\delta)\ne0$ \unsure{partt|partout} sur $S$)
% slanted
\margin{indépendant des résultats de Mumford}
\[
\text{``}\quad\varphi_L(\mathrm{cl}(L))=\chi\deleted{\unsure{\otimes}}L'^{-\frac1\chi}\quad\text{''}
\qquad\text{où }\chi=\chi(\delta)\in H^0(S,\mathbb Z).
\]
Au niveau du $\mathrm{ch}_B$, ceci s'écrit
\margin{$D'\in\mathrm{Gr}^1$ défini par $L'$}
\[
\begin{aligned}
\mathrm{ch}\,\varphi_L(\deleted{\unsure{\mathrm{cl}}}\unsure{\lambda|L|\mathrm{cl}(L)})&=\chi\exp\frac{-D'}{\chi}=\chi-D'+\frac1{2\chi}D'^{2}-\frac1{3!\chi^2}D'^{3}\ldots\\
&\quad\inserted{+(-1)^n\frac1{n!\chi^{n-1}}D'^{n}},
\end{aligned}
\]
et au niveau des classes de Chern elles-mêmes
\[
\begin{aligned}
c(\varphi_L(\unsure{\lambda|L|\mathrm{cl}(L)}))&=\left(1-\frac{D'}{\chi}\right)^\chi\\
&=1-D'+\frac{\chi(\chi-1)}{2\chi}D'^{2}-\frac{\chi(\chi-1)(\chi-2)}{3!\chi^2}D'^{3}+\ldots\\
&\quad+(-1)^n\frac{\chi(\chi-1)\ldots(\chi-n\unsure{+|-}1)}{n!\chi^{n-1}}D'^{n}.
\end{aligned}
\]

\margin{?}
Ces formules sont-elles bien vraies mod torsion ? Elles sont vraies au moins sur un corps alg clos, \uline{mod \uline{équivalence} cohomologique} $l$-adique (tout $l\ne\mathrm{car}.k$). Ceci donne des \margin{!} propriétés de divisibilité intéressantes de $\delta'$ savoir $\deleted{\unsure{\delta'}}\deleted{\unsure{\chi}}\frac{(\chi-1)\ldots(\chi-k+1)}{k!\,\chi^{k-2}}\,\delta'^{k}$ est ``entier'' i.e provient d'un cycle à coeff. entiers
% archive sheet number: 8
